\noindent\textbf{Lectura cuantitativa de los resultados.} En el texto en español ($S_o=36000$~B), Huffman alcanza $H(X)=4.4170$ y $L=4.4543$ bits/símbolo ($\eta=99.16$\,\%), reduciendo el archivo a $20166$~B (43.98\,\% de ahorro), mientras LZ77 aprovecha repeticiones de palabras y frases para llegar a $5131$~B (85.75\,\% de ahorro) y RLE lo expande a $107373$~B (-198.26\,\%). En la entrada homogénea de $32\,768$~B, RLE logra una razón $R=1560.38$ ($S_c=21$~B, 99.94\,\% de ahorro). En contraste, en bytes pseudoaleatorios y en el PNG ya comprimido ningún método logra ahorro: Huffman añade $+529$~B por tabla canónica y LZ77 añade $+4604$~B sobre el PNG por banderas de literales.\par
